% !TeX program = pdflatex
% papers/conversa_andar_peano.tex — Andar Peano no corpus (leve; depois da fundação).
\documentclass[11pt,a4paper]{article}
\usepackage[utf8]{inputenc}\usepackage[T1]{fontenc}\usepackage[brazil]{babel}
\usepackage[margin=2.4cm]{geometry}
\usepackage[hidelinks]{hyperref}
\newcommand{\code}[1]{\texttt{\small #1}}
% ─────────────────────────────────────────────────────────────────────────────
%  papers/gkcapa.tex — a capa do Reino, mínima, para os papers em `article`.
%
%  O estilo.tex completo é do `report` (títulos de capítulo, skak, …). Aqui fica
%  só o que a capa precisa, para o paper continuar a compilar sozinho com
%  pdflatex. O tradutor próprio (tests/tex) carrega o estilo.tex da raiz / o
%  symlink papers/estilo.tex e avalia o mesmo `\gkcapa`.
% ─────────────────────────────────────────────────────────────────────────────
\definecolor{ouro}{HTML}{B8912F}
\definecolor{ouroclaro}{HTML}{F3E9CE}
\definecolor{tinta}{HTML}{1E1B16}
\definecolor{regua}{HTML}{6E6858}
\providecommand{\gknota}{\fontsize{7.62}{11.05}\selectfont}
\providecommand{\gktexto}{\fontsize{10.50}{15.22}\selectfont}
\providecommand{\gksub}{\fontsize{12.33}{17.87}\selectfont}
\providecommand{\gksec}{\fontsize{14.47}{20.98}\selectfont}
\providecommand{\gkcap}{\fontsize{16.99}{24.63}\selectfont}
\providecommand{\gktit}{\fontsize{23.42}{33.95}\selectfont}
\providecommand{\Prj}{\mathbb{P}}
\providecommand{\Trf}{\mathcal{T}}
\providecommand{\gkcapa}[3]{%
\title{\vspace{-0.6cm}
{\color{ouro}\rule{\textwidth}{1.4pt}}\\[6mm]
\textbf{\gktit\color{tinta} Reino Dourado}\\[3mm]{\gksec\itshape\color{regua} Golden Kingdom}\\[8mm]
{\gkcap\scshape\color{ouro} #1}\\[5mm]
{\gksub\color{regua} #2}\\[2mm]
{\gktexto\color{regua}\itshape #3}\\[7mm]
{\color{ouro}\rule{\textwidth}{0.6pt}}\\[9mm]{\gknota\color{regua}\begin{minipage}{\textwidth}\setlength{\parindent}{0pt}%
\textbf{Aviso de Isenção Algorítmica:} O universo, as leis e as entidades contidas nestas páginas ---
incluindo felinos matriciais, aves de rapina algébricas e amuletos de controle de fluxo --- são
construções de pura ficção formal. Esta obra não descreve a realidade física, não serve como
engenharia reversa do cosmos e não constitui doutrina religiosa. Qualquer semelhança com bugs reais do seu
sistema operacional é mera coincidência (ou falta de reza). Prossiga por sua própria conta e
risco.\end{minipage}}}
\author{}\date{}}

\begin{document}
\gkcapa{Andar Peano}
       {sobe a torre depois das fundações}
       {retração, Maestro, Metrónomo --- sem aula de Lei 9}
\maketitle
\begin{abstract}
Pares curtos para o próximo andar da escada.
Fundação primeiro: \code{papers/torre\_fundacao.tex}.
Detalhe: \code{papers/corpo\_peano.tex}.
\end{abstract}

\section{o que é o corpo de peano}
Torre com retrações $\pi_k$ e $d_{k+1}=2d_k$.
Ver \code{papers/corpo\_peano.tex} --- depois de \code{papers/torre\_fundacao.tex}.

\section{o que faz o maestro}
Projecta: $\mathcal{P}_k=\mathrm{tick}\circ\mathrm{batuta}\circ\Pi$ realiza $\pi_k$.
Não compõe o corpo. \code{papers/corpo\_peano.tex}.

\section{o que faz o metronomo}
Lê a volta: $\lambda^{+}+\lambda^{-}=0$ atesta o fecho.
\code{papers/corpo\_peano.tex}.

\section{o que e histerese aqui}
Memória da borda sob a batuta $I\in\{-1,0,+1\}$ --- sem Lei~9.
\code{papers/corpo\_peano.tex}.

\section{como sobe a torre peano}
$A_0\xleftarrow{\pi_0}A_1\xleftarrow{\pi_1}\cdots$ com a dimensão a dobrar.
\code{papers/corpo\_peano.tex}.

\section{o que é a retracao pi\_k}
$\pi_k\colon A_{k+1}\to A_k$ com $\pi_k\circ\iota_k=\mathrm{id}$ --- desce um andar.
\code{papers/corpo\_peano.tex}.

\section{o que e controle de histerese}
Seleccionar variantes em $\mathcal{V}_k$ sem destruir o estado --- gerar, medir, seleccionar, reter.
\code{papers/corpo\_peano.tex}.

\section{mostra o corpo de peano}
\code{papers/corpo\_peano.tex} --- depois de \code{papers/torre\_fundacao.tex}.

\section{falemos de peano}
Andar da torre: retração, Maestro, Metrónomo. Fundação primeiro.
\code{papers/corpo\_peano.tex}.

\section{falemos do maestro}
Projecta o suporte; não inventa estrutura.
\code{papers/corpo\_peano.tex}.

\section{falemos do metronomo}
Atesta a volta (resíduo $\lambda^{+}+\lambda^{-}=0$).
\code{papers/corpo\_peano.tex}.

\section{o que e a batuta}
Canaliza $I\in\{-1,0,+1\}$ no Maestro --- sentido do MOVE.
\code{papers/corpo\_peano.tex}.

\section{como a torre dobra a dimensao}
$d_{k+1}=2d_k$ a cada andar.
\code{papers/corpo\_peano.tex}.

\section{o que e o par complexo do andar}
Par conjugado $(R,I)$ com $\Phi(x,y)=x+iy$ --- não o andar inteiro é $\mathbb{C}$.
\code{papers/corpo\_peano.tex} \S pares-complexos.

\section{o que e a leitura polar}
$(r,\theta)$: módulo = ocupação; fase = orientação; histerese = anel+setor.
\code{papers/corpo\_peano.tex}.

\section{a torre de i prova o peano?}
Não --- modelo de teste: $T_{k+1}=i^{T_k}$. Maestro = retração; $i^{T}$ = iteração.
\code{papers/corpo\_peano.tex}.

\section{o que e a formula de euler no peano}
No par: $\exp(tJ)=c(t)+s(t)J$ com $J^{2}=-1$; realização $J\mapsto i$ dá $e^{i\theta}=\cos\theta+i\sin\theta$.
\code{papers/corpo\_peano.tex} \S euler-peano.

\section{o que e lyapunov dualizado}
Três leituras: clássico $\lambda$; frente/volta $(\lambda^{+},\lambda^{-})$ com soma $0$; polar conjugado $|\mu_{\pm}|=|F'|$.
Não colar. \code{papers/corpo\_peano.tex} \S lyap-dual-releitura; \code{tests/lyapunov\_refletido.c}.

\section{o que e a rede neural dual}
Teorema: retenção na banda; realização dual sob conjugação $F_H=D F_P^{-1} D^{-1}$; $\lambda_P+\lambda_H=0$.
Hopfield = memória da volta, não inversão. Pesos $W=-I$ pela Lei 1.
\code{papers/corpo\_peano.tex} \S rede-dual.

\section{o que e a dualidade cruzada}
Imaginário controla o módulo seguinte; real controla a fase seguinte.
\code{papers/corpo\_peano.tex}.

\end{document}
